\documentclass[12pt,a4paper,oneside]{ctexart}
\usepackage{amsmath,amssymb,graphicx}
\usepackage[margin=2.5cm]{geometry}
\title{第二章 基本结构：集合、函数、序列、求和和矩阵}
\begin{document}

\section*{集合}

其他所有的离散结构都建立于集合的基础之上

Any collection of objects is called a set.

集合用于把对象组织在一起。

定义：集合是不同对象的一个无序的聚集，对象也称为集合的元素(element)或成员(member),集合包含(contain)它的元素。\\
我们使用 \(a\in A\) 来表示a是集合A中的一个元素，记号 \(a\notin A\) 表示a不是集合A中的一个元素。

\begin{quote}
通常用大写字母表示集合，而用小写字母表示集合中的元素。
\end{quote}

集合的描述方法：

\begin{itemize}
\item
  花名册方法\{a,b,c,d\},就是把所有元素穷举
\item
  用"...",比如\{...,-3,-1\},适用于那些元素之间规律十分明显的
\item
  使用集合构造器(set builder)符号\{x| P(x)\}
\item
  文氏图
\end{itemize}

\begin{quote}
值得注意的是，尽管集合常用来聚集具有共同性质的元素，但也不妨碍集合拥有表面上看起来毫不相干的元素，比如\{a, 2, Alice\}等等
\end{quote}

定义：两个集合相等当且仅当它们拥有相同的元素，如果A和B是集合，那么 \(\forall x(x\in A\leftrightarrow x\in B)\) ，记为A=B

注意：

\begin{itemize}
\item
  集合中元素的排列顺序无关紧要(无序性)
\item
  集合中元素出现不止一次也没关系，\{1,2\}与\{1,2,2\}是相同的集合(互异性)
\end{itemize}

空集：不含任何元素的集合，记为 \(\varnothing\) ，也可以记为\{ \}

单元素集：只含有1个元素的集合。

\begin{quote}
但是我们要区分 \(\varnothing\) 和\{ \(\varnothing\) \},前者为空集，而后者为单元素集。
\end{quote}

\subsection*{文氏图}

文氏图是用来描述集合的图表

在文氏图中全集(universal set)U,包含所考虑的全部对象，用矩形框来表示。\\
在矩形框内部，圆形或其他几何图形用于表示集合。有时候用点来表示集合中特定的元素。

文氏图常用于表示集合之间的关系。

\subsection*{子集}

定义：集合A是集合B的子集并且集合B是集合A的超集当且仅当A的每个元素也是B的元素，我们使用 \(A\subseteq B\) 表示集合A是集合B的子集。\\
另外，如果我们要强调集合B是集合A的超集，我们使用 \(B\supseteq A\) (当然这两种表达方式是一致的)

换句话说， \(A\subseteq B\) 当且仅当 \(\forall x(x\in A\rightarrow x\in B)\) 为真。

对于任意一个集合S，都有下面的定理：\\
定理：1） \(\varnothing\subseteq S\) 2） \(S\subseteq S\)

当我们要强调集合A是集合B的子集但是 \(A\neq B\) 时，我们就写成 \(A\subset B\) 并说A是B的真子集，也就是说B中有一个元素并不是A中的元素。

A是B的真子集当且仅当 \(\forall x(x\in A\rightarrow x\in B)\wedge\exists x(x\in B\wedge x\notin A)\)

但如果我们要证明两集合相等，我们需要证明 \(A\subseteq B\) 和 \(B\supseteq A\)

\subsection*{集合的大小}

定义：令S为集合。如果S中恰有n个不同元素(distinct elements)，这里n是非负整数，我们就说S为有限集，而n是S的基数(cardinality,也称为势)。S的基数记为| S|

\begin{quote}
基数作为衡量有限集合大小的一个重要常用语
\end{quote}

定义：一个集合是无限的，如果它不是有限的(呃。。。。)

\subsection*{幂集}

定义：给定集合S，S的幂集(power set)是集合S所有子集的集合。S的幂集记为P(S).

\begin{quote}
注意，该集合本身以及空集也是这个幂集中的元素。
\end{quote}

\begin{quote}
若| S|=n，则| P(S)|= \(2^n\) .
\end{quote}

\begin{quote}
\(x\in P(S)\Rightarrow x\subseteq S\)
\end{quote}

\begin{quote}
\(x\in S\Rightarrow \{x\}\subseteq P(S)\)
\end{quote}

\begin{quote}
\(S\in P(S)\)
\end{quote}

但我们要注意空集的幂集，

\begin{itemize}
\item
  P( \(\varnothing\) )=\{ \(\varnothing\) \},但是
\item
  P(\{ \(\varnothing\) \})=\{ \(\varnothing\) ,\{ \(\varnothing\) \}\},
\item
  P(\{ \(\varnothing\) ,\{ \(\varnothing\) \}\})=\{ \(\varnothing\) ,\{ \(\varnothing\) \},\{\{ \(\varnothing\) \}\},\{ \(\varnothing\) ,\{ \(\varnothing\) \}\}\}
\end{itemize}

\[\begin{gathered}
P(A)\in P(B)\Rightarrow P(A)\subseteq B\\
A\in P(A)\subseteq B\Rightarrow A\in B\\
\text{也就是说}，P(A)\in P(B)\Rightarrow A\in B
\end{gathered}\]

我们还有：

\[A\subseteq B\Rightarrow P(A)\subseteq P(B)\]

\subsection*{笛卡尔积}

由于集合是无序的，所以我们就需要用一种不同的结构来表示有序的聚集。这就是有序n元组。

定义：有序n元组(ordered n-tuple)( \(a_1,a_2,...,a_n\) )是以 \(a_1\) 为第一个元素， \(a_2\) 为第二个元素，...， \(a_n\) 为第n个元素的\textbf{有序}聚集。

两个有序n元组相等当且仅当每一对对应的元素都相等。

\begin{quote}
有序二元组也被称为序偶(ordered pairs)。
\end{quote}

定义：令A和B为集合。A和B的笛卡尔积(Cartesian product)用 \(A\times B\) 表示，是所有序偶(a,b)的集合，其中 \(a\in A\) 且 \(b\in B\) 。于是

\(A\times B=\{(a,b)|a\in A\wedge b\in B\}\)

注意：

\begin{itemize}
\item
  \(\varnothing\times A=A\times\varnothing=\varnothing\)
\item
  \(A\times B\) 与 \(B\times A\) 一般不相等，除非A= \(\varnothing\) 或 \(B=\varnothing\) 或 \(A=B\)
\end{itemize}

定义：集合 \(A_1,A_2,...,A_n\) 的笛卡尔积用 \(A_1\times A_2\times...\times A_n\) 来表示，是有序n元组 \((a_1,a_2,...,a_n)\) 的集合，其中 \(a_i\in A_i,i=1,2,...,n\)

\(A_1\times A_2\times...\times A_n\) =\{ \((a_1,a_2,...,a_n)\) | \(a_i\in A_i,i=1,2,...,n\) \}

\begin{itemize}
\item
  值得注意的是，当A，B，C是集合时， \((A\times B)\times C\) 与 \(A\times B\times C\) 是不同的，前者是有序二元组，而后者是有序三元组。
\end{itemize}

我们用 \(A^2\) 来表示 \(A\times A\) ,即集合A与自身的笛卡尔积，更一般地

\(A^n=\{(a_1,a_2,...,a_n)|a_i\in A,i=1,2,...,n\}\)

笛卡尔积 \(A\times B\) 的一个子集R被称为从集合A到集合B的关系。从集合A到自身的一个关系称为A上的一个关系。

\section*{集合运算}

实际上，命题演算和集合理论都是布尔代数这一代数系统的实例。集合理论上的操作，是依据与命题演算操作相对应而定义的。

通常我们会有一个全集，而所有集合都假定是全集的子集。

\subsection*{引言}

定义一：集合A与集合B的并集(Union)，用 \(A\cup B\) 来表示，仍是一个集合，包含A或B中或同时在A和B中的元素。

\[\begin{gathered}
A\cup B=\{x|x\in A \vee x\in B\}
\end{gathered}\]

\begin{itemize}
\item
  \(A\subseteq A\cup B\), \(B\subseteq A\cup B\)
\item
  \(A\subseteq C,B\subseteq C\Rightarrow A\cup B\subseteq C\)
\item
  \(|A\cup B|\leqslant|A|+|B|\)
\item
  \(A\cup B=B\Leftrightarrow A\subseteq B\)
\end{itemize}

定义二：集合A与集合B的交集(Intersection)，用 \(A\cap B\) 来表示，是一个集合，它包含同时在A和B中的元素。

\[\begin{gathered}
A\cap B=\{x|x\in A\wedge x\in B\}
\end{gathered}\]

\begin{itemize}
\item
  \(A\cap B\subseteq A\), \(A\cap B\subseteq B\)
\item
  \(C\subseteq A,C\subseteq B\Rightarrow C\subseteq A\cap B\)
\item
  \(|A\cap B|\leqslant|A|\), \(|A\cap B|\leqslant|B|\)
\item
  \(A\cap B= A\Leftrightarrow A\subseteq B\)
\end{itemize}

定义三：两个集合是不相交的，如果他们的交集为空集。

容斥原理： \(|A\cup B|=|A|+|B|-|A\cap B|\)

定义四：集合A和集合B的差集(Difference of A and B)

A-B=\{ \(x|x\in A\wedge x\notin B\) \}= \(A\cap\overline{B}\)

定义五：集合A的补集(the complement of a set),用 \(\overline{A}\) 来表示，其实就是U-A

定义六：集合A和集合B的对称差(Symmetric difference)

\(A\oplus B=(A\bigcup B)-(A\bigcap B)=(A-B)\bigcup(B-A)\)

\subsection*{集合恒等式}

\begin{itemize}
\item
  \(A\cap U=A\)
\item
  \(A\cup\varnothing=A\)(恒等律)
\item
  \(A\cup U=U\)
\item
  \(A\cap \varnothing=\varnothing\)(支配律)
\item
  \(A\cup A=A\)
\item
  \(A\cap A=A\)(幂等律)
\item
  \(\overline{(\overline{A})}=A\)(补律)
\item
  \(A\cup B=B\cup A\)
\item
  \(A\cap B=B\cap A\)(交换律)
\item
  \(A\cup(B\cup C)=(A\cup B)\cup C\)
\item
  \(A\cap(B\cap C)=(A\cap B)\cap C\)(结合律)
\item
  \(A\cup(B\cap C)=(A\cup B)\cap(A\cup C)\)
\item
  \(A\cap(B\cup C)=(A\cap B)\cup(A\cap C)\)(分配律)
\item
  \(\overline{A\cap B}=\overline{A}\cup\overline{B}\)
\item
  \(\overline{A\cup B}=\overline{A}\cap\overline{B}\)(德摩根律)
\item
  \(A\cup(A\cap B)=A\)
\item
  \(A\cap(A\cup B)=A\)(吸收律)
\item
  \(A\cup\overline{A}=U\)
\item
  \(A\cap\overline{A}=\varnothing\)(互补律)
\end{itemize}

\section*{函数}

\subsection*{引言}

定义：令A和B为非空集合。从A到B的函数f是对元素的一种指派，对A的每个元素恰好指派B的一个元素。如果B中元素b是唯一由函数f指派给A中元素a的，则我们就写成f(a)=b。如果f是从A到B的函数，就记为 \(A\rightarrow B\) 。

函数f: \(A\rightarrow B\) 也能从A到B的关系( \(A\times B\) 的子集)来定义(用序偶(a,b)去定义f(a)=b)。

\begin{quote}
函数有时也被称为映射(mapping)或者变换(transformation)。
\end{quote}

定义：如果f是从A到B的函数，我们说A是f的定义域(domain),而B是f的陪域(codomain) 。如果f(a)=b，我们说b是a的像(image)，而a是b的原像(preimage)。f的值域(range)或像是A中元素所有像的集合。如果f是从A到B的函数，我们说f把A映射到B。

\begin{quote}
注意区分陪域与值域，实际上值域是陪域的子集。
\end{quote}

定义函数的时候，我们需要指定它的定义域、陪域、定义域中元素到陪域的映射。当两个函数的这三个因素相同时，我们称这两个函数是相等的。

如果一个函数的陪域是实数集合，那么该函数称为实值函数。如果其陪域为整数集合，那么该函数称为整数值函数。具有相同定义域的两个实值函数或两个整数值函数可以相加和相乘。

定义：令 \(f_1,f_2\) 是从A到R的函数，那么 \(f_1+f_2,f_1f_2\) 也是从A到R的函数

\[\begin{gathered}
\forall x\in A\\
(f_1+f_2)(x)=f_1(x)+f_2(x)\\
(f_1f_2)(x)=f_1(x)f_2(x)
\end{gathered}\]

\subsection*{一对一函数和映上函数}

定义：若对于函数f的定义域中所有的a和b都有f(a)=f(b)蕴含a=b，则称该函数f(x)为一对一(one-to-one)函数或者单射(injection)函数。

换句话说，像只对应一个原像。

\(\forall x\forall y(f(x)=f(y)\rightarrow x=y)\)

\(\forall x\forall y(x\neq y\rightarrow f(x)\neq f(y))\)

定义：递增/严格递增/递减/严格递减/

定义：若对函数f的陪域中任意元素b都有元素a使得f(a)=b，则称该函数f为映上(onto)函数或满射(surjection)函数

定义：既是单射函数又是满射函数的函数为双射(bijection)函数或一一对应(one-to-one correspondence)函数。此时两集合具有相同的基数(势)。

假设f: \(A\rightarrow B\)

要证明f是单射的： \(\forall x,y\in A,\text{如果}f(x)=f(y),\text{则}x=y\)

要证明f不是单射的：找到特定的 \(x,y\in A\) 使得 \(x\neq y\) 且 \(f(x)= f(y)\)

要证明f是满射的：考虑 \(\forall y\in B\) ,并找到一个元素 \(x\in A\) 使得f(x)=y

要证明f不是满射的：找到一个特定的 \(y\in B\) ,使得对于任意 \(x\in A\) 有f(x) \(\neq y\)

\subsection*{反函数和函数合成}

定义：令f为从集合A到集合B的双射函数。f的反函数指派给B中的元素b是A中使得f(a)=b唯一元素a。f的反函数用 \(f^{-1}\) 来表示。

只有双射函数才能有逆函数(也被称为是可逆的)。

定义：(函数的合成)令g为集合A到集合B的函数，f是从集合B到集合C的函数，函数f和g的合成(composition)记作f \(\circ\) g。定义为对任意A中的元素a，有(f \(\circ\) g)(a)=f(g(a))。

\subsection*{函数的图}

函数f的图像其实是序偶集合(笛卡尔积 \(A\times B\) 的子集)

\(\{(a,b)|a\in A\ and\ f(a)=b\}\)

\subsection*{一些重要的函数}

一、低函数：

低函数指派给实数x的是小于或等于x的最大整数，用 \(\lfloor x\rfloor\) 来表示(其实就是下取整函数(floor))。

二、顶函数

顶函数指派给实数x的是大于或等于x的最小整数，用 \(\lceil x\rceil\) (其实就是上取整(ceiling)函数)

\begin{itemize}
\item
  \(x-1<\lfloor x\rfloor\leqslant x\leqslant\lceil x\rceil x+1\)
\item
  \(\lfloor -x\rfloor=-\lceil x\rceil\)
\item
  \(\lceil -x\rceil=-\lfloor x\rfloor\)
\item
  \(\lfloor x+m\rfloor=\lfloor x\rfloor+m\) 当m是整数时
\item
  \(\lceil x+m\rceil=\lceil x\rceil+m\) 当m是整数时
\end{itemize}

\section*{集合的基数}

\subsection*{引言}

定义：集合A和集合B有相同的基数(cardinality)，当且仅当存在从A到B的一个一一对应。当A和B有相同的基数时，就写成| A|=| B|。

对于无限集，基数的定义提供了一个衡量两个集合相对大小的方法，而不是衡量一个集合大小的办法。

例：设区间(a,b)中含有实数的集合记为集合A，区间(0,1)中含有的实数集合记为集合B，证明集合A,B具有相同的基数。

\[\begin{gathered}
\text{令}f(x)=\frac{x-a}{b-a},x\in(a,b)\\
\text{显然}f(x)\text{是一个从集合}(a,b)\text{到集合}(0,1)\text{的双射函数}\\
\text{故两集合具有相同的基数}
\end{gathered}\]

定义：如果存在一个从A到B的一对一函数，则A的基数小于或等于B的基数，并写成| A| \(\leqslant\) | B|。再者，当| A| \(\leqslant\) | B| 并且A和B有不同的基数时，我们说A的基数小于B的基数。

\subsection*{可数集合}

定义：一个集合或者是有限集或者与自然数集具有相同的基数，这个集合就称为可数的。一个集合不是可数的，就称为不可数的。(如果该集合与自然数集具有相同的基数，我们称其为无限可数的)

可数分为两种情况，一是有限可数，二是无限可数。

如果我们要证明一个集合是无限可数的，那么我们必须给出一个该集合与正整数集合之间的一一映射，在这之中，该集合可以做定义域，也可以做陪域。

我们记自然数集合的势为 \(\aleph_0\) 。

一些无穷可数的集合：

\begin{itemize}
\item
  整数集合Z
\item
  正偶数集合E
\item
  正有理数集合 \(Q^+\)
\end{itemize}

1.整数集合Z

\[\begin{gathered}
Z=\{0,-1,1,-2,2,...\}\\
f(i)=
\begin{cases}
2|i|\;\;\;i<0\\
1\;\;\;i=0\\
2|i|+1\;\;\;i>0\\
\end{cases}\\
f\text{是一个从集合}Z\text{到集合}N\text{的双射函数}\\
\text{故}Z\text{是一个无限可数的集合}
\end{gathered}\]

2.正偶数集合E

\[\begin{gathered}
E=\{2n|n\in N\}\\
f(x)=2x\;\;\;x\in N\\
f\text{是一个从集合}N\text{到集合}E\text{的一个双射}\\
\text{故集合}E\text{是一个无限可数的集合}\\
\text{同理可得正奇数集合也是一个无限可数集合}
\end{gathered}\]

\begin{itemize}
\item
  值得注意的是，尽管正偶数集合是自然数集合的真子集，但是正偶数集合与自然数集合的势是相同的。
\end{itemize}

3.正有理数集合 \(Q^+\)

\[\begin{gathered}
\forall x\in Q^+,x=\frac{q}{p}\;\;p,q\in N\text{且互质}\\
\text{令集合}S=\{(p,q)|p,q\in N\}=N\times N\\
\text{令函数}f(x)=f(\frac{q}{p})=(p,q)\;\;x\in Q^+\\
\text{显然}f\text{是一个从}Q^+\text{到}S\text{的单射函数},\text{故}|Q^+|\leqslant|S|\\
\begin{matrix}
(1,1)&(1,2)&(1,3)&...&(1,q)&...\\
(2,1)&(2,2)&(2,3)&...&(2,q)&...\\
(3,1)&(3,2)&(3,3)&...&(3,q)&...\\
...\\
(p,1)&(p,2)&(p,3)&...&(p,q)&...\\
...
\end{matrix}\\
\text{对于上面的矩阵}，\text{我们以对角线顺序向自然数集对应}\\
(1,1)\text{对应}1，(1,2)\text{对应}2，(2,1)\text{对应}3...\text{以此类推}\\
\text{令}g(x)=g((p,q))=\frac{(p+q-2)(p+q-1)}{2}+p\;\;\;x\in S\\
g\text{是一个从}S\text{到}N\text{的满射函数},\text{故}|S|=|N|\\
\text{又因为}N\subseteq Q^+,\text{故}|N|\leqslant|Q^+|\\
\text{综上有}|N|=|Q^+|
\end{gathered}\]

\begin{itemize}
\item
  有理数集合Q也是无限可数的。
\end{itemize}

无限可数集合的性质：

\begin{itemize}
\item
  不存在一个无限集合比一个可数集合的基数还小
\item
  两个可数集合的并还是可数的。
\item
  有限个可数集合的并还是可数的。
\item
  可数的可数集合的并还是可数的。
\end{itemize}

康托尔对角化论证：

定理：在0和1之间的实数集合是不可数的。

\[\begin{gathered}
A=\{x|x\in(0,1)\wedge x\in R\}\\
B=\{\frac{1}{n+1}|n\in N\}\\
\text{显然}|N|=|B|，B\subseteq A\\
\text{所以}|N|\leqslant|A|\\
\text{假设集合}A\text{是无限可数的}，A=\{r_1,r_2,...,r_n,...\}\\
\text{将}A\text{中的每个数都已十进制的形式表示出来}：\\
\text{比如}\frac{1}{3}=0.3333333...,\frac{1}{2}=0.50000...\\
r_1=0.d_{11}d_{12}d_{13}d_{14}d_{15}d_{16}d_{17}...\\
r_2=0.d_{21}d_{22}d_{23}d_{24}d_{25}d_{26}d_{27}...\\
r_3=0.d_{31}d_{32}d_{33}d_{34}d_{35}d_{36}d_{37}...\\
...\\
\text{现在我们构造一个实数}x=0.x_1x_2x_3x_4x_5x_6...\\
\text{如果}d_{ii}\neq 3,x_i=3\\
\text{如果}d_{ii}=3,x_i=4\\
\text{我们发现}x\text{不同于集合}A\text{中的任意一个元素}\\
\text{但}x\text{确实属于区间}(0,1)\text{之间}，\text{说明假设错误}\\
\text{即集合}A\text{是无限不可数的}，\text{并且}|N|<|A|\\
\end{gathered}\]

定理：从负无穷大到正无穷大区间的实数个数在0和1之间的实数个数的基数是相同的。

\begin{quote}
提示：使用f(x)= \(\frac{2tanx}\pi\) 来证明。
\end{quote}

例：证明集合{[}0,1{]}的势为 \(\aleph\) 。

\[\begin{gathered}
A=[0,1],B=(0,1)\\
\text{故}B\subseteq A,\text{则有}|B|\leqslant|A|\\
\text{令}g(x)=\frac{1}{2}x+\frac{1}{4}\;\;\;x\in [0,1]\\
g\text{是一个从}[0,1]\text{到}[\frac{1}{4},\frac{3}{4}]\text{的双射函数}\\
\text{所以}|A|\leqslant|B|\\
\text{综上}|A|=|B|
\end{gathered}\]

我们记实数集合R的势为 \(\aleph\) 。

连续性假设：不存在一个基数a，使得 \(\aleph_0<a<\aleph\)

\end{document}
